% Copyright (c) 2026 Twin Vector Labs LLC.
% All rights reserved.
%
% Pedagogical study document: from basic quantum spin to the entangled
% readout of the emergent proton spin, as implemented in the Tessera codebase.
\documentclass[11pt,oneside]{article}

\usepackage[margin=1in]{geometry}
\usepackage{amsmath}
\usepackage{amssymb}
\usepackage{mathtools}
\usepackage{braket}
\usepackage{bm}
\usepackage{enumitem}
\usepackage{booktabs}
\usepackage{array}
\usepackage{xcolor}
\usepackage[colorlinks=true,linkcolor=blue!55!black,citecolor=green!45!black,%
            urlcolor=blue!55!black]{hyperref}

% --- convenience macros -------------------------------------------------------
\newcommand{\R}{\mathbb{R}}
\newcommand{\C}{\mathbb{C}}
\newcommand{\Z}{\mathbb{Z}}
\newcommand{\halff}{\tfrac{1}{2}}
\newcommand{\id}{\mathbb{1}}
\newcommand{\Sig}{\Sigma}
\newcommand{\Tr}{\operatorname{Tr}}
\newcommand{\eqdef}{\vcentcolon=}
\DeclareMathOperator{\sgn}{sgn}
% physics.sty is unavailable on this TeX install; provide the few macros we use.
\newcommand{\expval}[1]{\langle #1 \rangle}
\newcommand{\mel}[3]{\langle #1 | #2 | #3 \rangle}
\newcommand{\norm}[1]{\left\lVert #1 \right\rVert}
\newcommand{\code}[1]{\texttt{#1}}
\newcommand{\up}{\uparrow}
\newcommand{\down}{\downarrow}

\title{\bfseries From Spin--\textonehalf{} to the Entangled Readout of the Proton Spin\\[4pt]
\large A study companion to the Tessera emergent--proton construction}
\author{Tessera theory notes \quad(\,part of issue \#410\,)}
\date{\today}

\begin{document}
\maketitle

\begin{abstract}
\noindent
This is a self--contained study document. It walks, rung by rung, from the
quantum mechanics of a single spin--\textonehalf{} particle up to the central
problem solved (and partly still open) in the Tessera codebase: reading the
\emph{composite} total spin of the emergent proton bound state and recovering
the value $J^2=\tfrac34$. Every symbol is defined where it first appears. The
physics ladder (Sections~\ref{sec:onespin}--\ref{sec:identical}) is standard
and can be checked against any quantum--mechanics or particle--physics
textbook; the geometry and Dirac--K\"ahler material
(Sections~\ref{sec:lattice}--\ref{sec:dk}) sets up the discrete language the
code speaks; the last three sections
(Sections~\ref{sec:emergent}--\ref{sec:summary}) connect that language to the
actual modules and findings. The pivotal conceptual point, established in
Section~\ref{sec:threespin} and paid off in
Sections~\ref{sec:composite}--\ref{sec:totalspace}, is that the proton's spin
is an \emph{entangled} three--body state, and that any readout which reduces
each quark to its own qubit \emph{provably} cannot reproduce $\tfrac34$.
\end{abstract}

\tableofcontents
\bigskip

% =============================================================================
\section{One spin--\textonehalf{}: the smallest quantum top}
\label{sec:onespin}
% =============================================================================

A spin--\textonehalf{} particle is the simplest non--trivial quantum system: a
two--dimensional complex Hilbert space $\mathcal H = \C^2$. We pick an
orthonormal basis and name the basis vectors
\begin{equation}
  \ket{\up} = \begin{pmatrix} 1 \\ 0 \end{pmatrix}, \qquad
  \ket{\down} = \begin{pmatrix} 0 \\ 1 \end{pmatrix},
\end{equation}
``spin up'' and ``spin down'' along a chosen quantization axis (conventionally
$\hat z$). A general (pure) state is a normalized complex superposition
$\ket{\psi} = \alpha\ket{\up} + \beta\ket{\down}$ with
$|\alpha|^2+|\beta|^2 = 1$.

\paragraph{The Pauli matrices.}
The three Hermitian, traceless $2\times2$ Pauli matrices are
\begin{equation}
  \sigma_1 = \begin{pmatrix} 0 & 1 \\ 1 & 0 \end{pmatrix}, \quad
  \sigma_2 = \begin{pmatrix} 0 & -i \\ i & 0 \end{pmatrix}, \quad
  \sigma_3 = \begin{pmatrix} 1 & 0 \\ 0 & -1 \end{pmatrix}.
\end{equation}
They satisfy
\begin{equation}
  \sigma_a\sigma_b = \delta_{ab}\,\id + i\,\epsilon_{abc}\,\sigma_c ,
  \label{eq:pauliprod}
\end{equation}
where $\delta_{ab}$ is the Kronecker delta, $\epsilon_{abc}$ is the totally
antisymmetric Levi--Civita symbol ($\epsilon_{123}=+1$), $\id$ is the
$2\times2$ identity, and repeated indices are summed (Einstein convention).
Equation~\eqref{eq:pauliprod} packages two facts: the \emph{anticommutator}
$\{\sigma_a,\sigma_b\} = 2\delta_{ab}\id$ (so each $\sigma_a^2=\id$), and the
\emph{commutator} $[\sigma_a,\sigma_b] = 2i\,\epsilon_{abc}\sigma_c$.

\paragraph{The spin operators.}
The physical spin observables (in units $\hbar=1$, used throughout) are
\begin{equation}
  S_a \eqdef \halff\,\sigma_a, \qquad a\in\{1,2,3\}.
\end{equation}
From $[\sigma_a,\sigma_b]=2i\epsilon_{abc}\sigma_c$ we get the defining
angular--momentum algebra
\begin{equation}
  [S_a, S_b] = i\,\epsilon_{abc}\,S_c .
  \label{eq:su2}
\end{equation}
This is the Lie algebra $\mathfrak{su}(2)$; any set of three operators obeying
\eqref{eq:su2} is a representation of spin/angular momentum, a fact we will use
repeatedly.

\paragraph{The Casimir.}
The total--spin--squared operator
\begin{equation}
  S^2 \eqdef S_1^2 + S_2^2 + S_3^2
\end{equation}
commutes with every $S_a$ (it is the \emph{Casimir} of $\mathfrak{su}(2)$).
Using $\sigma_a^2 = \id$,
\begin{equation}
  S^2 = \tfrac14\big(\sigma_1^2+\sigma_2^2+\sigma_3^2\big)
      = \tfrac14\cdot 3\,\id = \tfrac34\,\id .
  \label{eq:onecasimir}
\end{equation}
Comparing with the general eigenvalue $s(s+1)$ of $S^2$ gives
$s(s+1)=\tfrac34 \Rightarrow s=\tfrac12$: every state of $\C^2$ is a
spin--\textonehalf{} state, and $\tfrac34$ is the fingerprint of a single
spin--\textonehalf{} degree of freedom. \textbf{Keep the number $\tfrac34$ in
view}: it is the value the proton must also reproduce, for a quite different
(composite) reason.

\paragraph{Measurement and expectation values.}
A measurement of $S_a$ yields one of its eigenvalues $\pm\tfrac12$, and in the
state $\ket\psi$ the average (expectation value) over many measurements is
\begin{equation}
  \expval{S_a} \eqdef \mel{\psi}{S_a}{\psi} .
\end{equation}
The triple $\expval{\bm S}=(\expval{S_1},\expval{S_2},\expval{S_3})$ is the
\emph{Bloch vector}; for a pure state it has length
$|\expval{\bm S}|=\tfrac12$. Equivalently the state's density matrix is
$\rho = \halff(\id + 2\,\expval{\bm S}\cdot\bm\sigma)$. This Bloch picture is
exactly how the code will summarize the spin extracted at each ``hole'' in
Section~\ref{sec:emergent}, and the fact that a Bloch vector keeps \emph{only}
$\expval{\bm S}$ (three real numbers) is precisely what loses the entanglement
in Section~\ref{sec:composite}.

% =============================================================================
\section{Two spins: the singlet/triplet split $2\otimes2 = 1\oplus3$}
\label{sec:twospin}
% =============================================================================

\paragraph{Tensor products.}
Two particles live in the tensor--product space
$\mathcal H = \C^2\otimes\C^2 \cong \C^4$, spanned by
$\ket{\up\up},\ket{\up\down},\ket{\down\up},\ket{\down\down}$, where
$\ket{ab}\eqdef\ket a\otimes\ket b$. An operator acting on particle~1 alone is
$A\otimes\id$; on particle~2 alone, $\id\otimes B$. The \emph{total} spin
operator is the sum
\begin{equation}
  \bm S_{\mathrm{tot}} = \bm S^{(1)}\otimes\id + \id\otimes\bm S^{(2)},
\end{equation}
which still obeys \eqref{eq:su2} because the two single--particle algebras
commute with each other.

\paragraph{Total Casimir.}
With $\bm S^{(1)}\cdot\bm S^{(2)} \eqdef \sum_a S^{(1)}_a S^{(2)}_a$,
\begin{equation}
  S_{\mathrm{tot}}^2
  = (\bm S^{(1)})^2 + (\bm S^{(2)})^2 + 2\,\bm S^{(1)}\!\cdot\bm S^{(2)}
  = \tfrac34 + \tfrac34 + 2\,\bm S^{(1)}\!\cdot\bm S^{(2)}.
  \label{eq:twocasimir}
\end{equation}
The cross term $\bm S^{(1)}\!\cdot\bm S^{(2)}$ is what couples the two spins;
its eigenvalues sort the four states into multiplets.

\paragraph{Clebsch--Gordan: the decomposition.}
The representation--theory statement is
\begin{equation}
  \boxed{\,2\otimes 2 = 1 \oplus 3\,}
\end{equation}
i.e.\ $\C^2\otimes\C^2$ splits into a one--dimensional \emph{singlet} ($S=0$)
and a three--dimensional \emph{triplet} ($S=1$). Explicitly, in terms of the
total magnetic quantum number $m$ (the $S^z_{\mathrm{tot}}$ eigenvalue):
\begin{align}
  \text{triplet } (S=1):\quad
    & \ket{1,+1} = \ket{\up\up}, \notag\\
    & \ket{1,\phantom{+}0} = \tfrac{1}{\sqrt2}\big(\ket{\up\down}+\ket{\down\up}\big),
        \label{eq:triplet}\\
    & \ket{1,-1} = \ket{\down\down}; \notag\\[4pt]
  \text{singlet } (S=0):\quad
    & \ket{0,0} = \tfrac{1}{\sqrt2}\big(\ket{\up\down}-\ket{\down\up}\big).
        \label{eq:singlet}
\end{align}
The numerical prefactors $\pm\tfrac{1}{\sqrt2}$ are the \emph{Clebsch--Gordan
coefficients} for $\tfrac12\otimes\tfrac12$. Note the symmetry under exchange of
the two particles: the triplet is \emph{symmetric}, the singlet is
\emph{antisymmetric}. This symmetry bookkeeping is the seed of the Pauli
argument in Section~\ref{sec:identical}.

\paragraph{Eigenvalues.}
On the triplet, $S_{\mathrm{tot}}^2 = 1(1+1) = 2$; on the singlet,
$S_{\mathrm{tot}}^2 = 0(0+1) = 0$. Cross--checking with
\eqref{eq:twocasimir}: $\bm S^{(1)}\!\cdot\bm S^{(2)} = +\tfrac14$ on the
triplet and $-\tfrac34$ on the singlet, so
$S^2_{\mathrm{tot}} = \tfrac32 + 2(\tfrac14) = 2$ and
$\tfrac32 + 2(-\tfrac34) = 0$ respectively. Already here we see that
$S^2_{\mathrm{tot}}$ is \emph{not} a property of either particle alone: it is
fixed by the relative phase/sign in \eqref{eq:triplet}--\eqref{eq:singlet}.

% =============================================================================
\section{Three spins: $2\otimes2\otimes2 = 4\oplus2\oplus2$ and the entangled
proton}
\label{sec:threespin}
% =============================================================================

Now combine three spin--\textonehalf{} particles, the case directly relevant to
a baryon (three quarks). The Hilbert space is
$\C^2\otimes\C^2\otimes\C^2\cong\C^8$, with total spin operator
$\bm J \eqdef \bm S^{(1)}+\bm S^{(2)}+\bm S^{(3)}$ (we switch to $\bm J$ for the
three--body total, to match the baryon literature). We use $\ket{abc}$ for
$\ket a\otimes\ket b\otimes\ket c$, and write e.g.\ $\ket{uud}$ when we later
attach flavor labels $u,d$; here read $u\equiv\up$, $d\equiv\down$ as spin
states.

\paragraph{The decomposition.}
Building up $2\otimes2\otimes2 = (1\oplus3)\otimes2 = 2 \oplus (4\oplus 2)$,
\begin{equation}
  \boxed{\,2\otimes2\otimes2 = 4 \oplus 2 \oplus 2\,}
\end{equation}
a four--dimensional \emph{quartet} with $J=\tfrac32$ and \emph{two}
two--dimensional \emph{doublets} with $J=\tfrac12$. The corresponding Casimir
$J^2 = j(j+1)$ values are
\begin{equation}
  J=\tfrac32:\quad J^2 = \tfrac32\cdot\tfrac52 = \tfrac{15}{4}, \qquad\qquad
  J=\tfrac12:\quad J^2 = \tfrac12\cdot\tfrac32 = \tfrac34 .
\end{equation}
The quartet ($\tfrac{15}{4}$) is the $\Delta$ baryon's spin; the doublets
($\tfrac34$) are the proton/neutron's. The two doublets are
distinguished by their behaviour under permutations of the three particles:
they are \emph{mixed--symmetry} representations (neither fully symmetric nor
fully antisymmetric), conventionally split into a piece symmetric under
$1\!\leftrightarrow\!2$ and a piece antisymmetric under $1\!\leftrightarrow\!2$.
The quartet is \emph{fully symmetric}.

\paragraph{The fully--symmetric quartet.}
The top of the $J=\tfrac32$ tower is $\ket{\tfrac32,\tfrac32}=\ket{\up\up\up}$
($=\ket{uuu}$ with all spins up). It is manifestly symmetric under any exchange
of the three particles, and one checks directly $J^2\ket{uuu}=\tfrac{15}{4}\ket{uuu}$.

\paragraph{The mixed--symmetry $J=\tfrac12$ ``proton'' state.}
The piece of $J=\tfrac12,\ m=+\tfrac12$ that is \emph{symmetric} under
exchanging particles $1\!\leftrightarrow\!2$ is, up to normalization,
\begin{equation}
  \ket{\chi^{\,\mathrm{MS}}}
   \;\propto\; 2\ket{\up\up\down} - \ket{\up\down\up} - \ket{\down\up\up}
   \;=\; 2\ket{uud} - \ket{udu} - \ket{duu}.
  \label{eq:protonstate}
\end{equation}
(Here the labels $u,d$ are spin up/down; the same combinatorial object appears
again in Section~\ref{sec:identical} with $u,d$ reading as \emph{flavor}.) The
normalized state is $\ket{\chi}=\tfrac{1}{\sqrt6}\big(2\ket{uud}-\ket{udu}-\ket{duu}\big)$.

Let us \textbf{verify} $J^2\ket\chi = \tfrac34\ket\chi$. Write
$J^2 = \sum_{i}(\bm S^{(i)})^2 + 2\sum_{i<j}\bm S^{(i)}\!\cdot\bm S^{(j)}
     = \tfrac94\,\id + 2\big(\bm S^{(1)}\!\cdot\bm S^{(2)}
       + \bm S^{(1)}\!\cdot\bm S^{(3)} + \bm S^{(2)}\!\cdot\bm S^{(3)}\big)$.
The convenient identity for any two spin--\textonehalf's is
\begin{equation}
  \bm S^{(i)}\!\cdot\bm S^{(j)} = \halff\,P_{ij} - \tfrac14\,\id,
  \label{eq:swap}
\end{equation}
where $P_{ij}$ is the \emph{swap} operator exchanging particles $i$ and $j$
($P_{ij}\ket{\dots a_i\dots a_j\dots}=\ket{\dots a_j\dots a_i\dots}$). Summing
\eqref{eq:swap} over the three pairs,
\begin{equation}
  J^2 = \tfrac94\id + 2\Big[\halff(P_{12}+P_{13}+P_{23}) - \tfrac34\id\Big]
      = \tfrac34\id + (P_{12}+P_{13}+P_{23}).
  \label{eq:j2swap}
\end{equation}
So we only need the action of the three swaps on $\ket\chi$. By construction
\eqref{eq:protonstate} is symmetric in $1\!\leftrightarrow\!2$, so
$P_{12}\ket\chi=+\ket\chi$. For the others, act on the unnormalized
$\ket v = 2\ket{uud}-\ket{udu}-\ket{duu}$:
\begin{align}
  P_{13}\ket v &= 2\ket{duu} - \ket{udu} - \ket{uud}, \\
  P_{23}\ket v &= 2\ket{udu} - \ket{uud} - \ket{duu}.
\end{align}
Adding all three swap results,
\begin{equation}
  (P_{12}+P_{13}+P_{23})\ket v
    = \underbrace{(2-1-1)}_{=0}\ket{uud}
    + \underbrace{(-1-1+2)}_{=0}\ket{udu}
    + \underbrace{(-1+2-1)}_{=0}\ket{duu} = 0.
\end{equation}
Hence $(P_{12}+P_{13}+P_{23})\ket\chi = 0$ and, from \eqref{eq:j2swap},
\begin{equation}
  \boxed{\,J^2\ket\chi = \tfrac34\ket\chi.\,}
\end{equation}
The state \eqref{eq:protonstate} is a genuine $J=\tfrac12$ eigenstate. Notice it
is an \emph{entangled} superposition of three product kets with non--trivial
relative coefficients $(2,-1,-1)$ --- it cannot be written as a single
$\ket a\otimes\ket b\otimes\ket c$.

\paragraph{The crucial counterexample: a plain product is \emph{not} $\tfrac34$.}
Contrast \eqref{eq:protonstate} with the bare product $\ket{uud}=\ket{\up\up\down}$.
Using \eqref{eq:j2swap}: $P_{12}\ket{uud}=\ket{uud}$ (the first two are both
$\up$), while $P_{13}\ket{uud}=\ket{duu}$ and $P_{23}\ket{uud}=\ket{udu}$ are
\emph{different} kets. Therefore
\begin{equation}
  J^2\ket{uud} = \tfrac34\ket{uud} + \ket{uud} + \ket{duu} + \ket{udu},
\end{equation}
and the \emph{expectation value} (the product ket is normalized,
$\braket{uud|uud}=1$, and is orthogonal to $\ket{duu},\ket{udu}$) is
\begin{equation}
  \boxed{\,\expval{J^2}_{uud} = \mel{uud}{J^2}{uud} = \tfrac34 + 1 = \tfrac74.\,}
  \label{eq:product74}
\end{equation}
A bare product of three spins gives $\tfrac74$, \emph{not} $\tfrac34$. In fact
$\ket{uud} = \tfrac{1}{\sqrt6}\,(\text{the }J{=}\tfrac32\text{ part}) + \dots$
is a mixture of $J=\tfrac32$ and $J=\tfrac12$, and $\tfrac74$ is the
mixture--weighted average. \emph{This is the single most important conceptual
point of the whole document}: the proton's $\tfrac34$ lives in the relative
$(2,-1,-1)$ phases --- it is an \emph{entanglement} property of the three--body
state. Any procedure that treats the three quarks as independent (a product, or
equivalently a per--particle Bloch reduction) cannot land on $\tfrac34$; it
floors at the mixture. We will meet exactly this floor in the lattice readout of
Section~\ref{sec:composite}.

% =============================================================================
\section{Identical particles, exchange symmetry, and the baryon wavefunction}
\label{sec:identical}
% =============================================================================

Why does Nature pick the specific entangled combination \eqref{eq:protonstate}
for the proton, rather than $\ket{uud}$ or the quartet? The answer is the Pauli
principle applied to the quarks' \emph{total} wavefunction.

\paragraph{Fermions and antisymmetry.}
Quarks are spin--\textonehalf{} fermions. The total wavefunction of a system of
identical fermions must be \emph{totally antisymmetric} under the simultaneous
exchange of all labels of any two of them. A quark carries four kinds of label:
\begin{equation}
  \Psi = \psi_{\text{space}} \otimes \chi_{\text{spin}} \otimes
         \phi_{\text{flavor}} \otimes \theta_{\text{color}},
\end{equation}
and exchange acts on all four factors at once. Antisymmetry of the product means
the symmetries of the factors must \emph{multiply} to ``antisymmetric.''

\paragraph{A light touch of Young tableaux.}
The permutation group $S_3$ on three objects has three irreducible symmetry
types, conveniently labelled by Young diagrams (written here as partitions):
\[
  \text{symmetric} \;\sim\; [3], \qquad
  \text{mixed (MS,\,MA)} \;\sim\; [2,1], \qquad
  \text{antisymmetric} \;\sim\; [1,1,1].
\]
(Pictorially: $[3]$ is a single row of three boxes, $[1,1,1]$ a single column
of three, and $[2,1]$ the ``L''/hook shape.) In words: a single row =
totally symmetric; a single column = totally antisymmetric; the ``L''/hook
shape = the two--dimensional \emph{mixed--symmetry} representation, which comes
as a pair MS, MA --- symmetric resp.\ antisymmetric under the
$1\!\leftrightarrow\!2$ exchange.) The combination rule we need is the
\emph{Clebsch--Gordan of $S_3$}: the product of two mixed--symmetry doublets
contains exactly one totally symmetric piece,
\begin{equation}
  \text{MS}\otimes\text{MS} \;\supset\; \text{symmetric},
  \label{eq:msmssym}
\end{equation}
formed as $\tfrac{1}{\sqrt2}(\phi^{\mathrm{MS}}\chi^{\mathrm{MS}}
+ \phi^{\mathrm{MA}}\chi^{\mathrm{MA}})$.

\paragraph{Color is the antisymmetric singlet.}
Each quark carries an $SU(3)_{\text{color}}$ index $a\in\{r,g,b\}$. A baryon is
a color singlet, and the unique antisymmetric singlet of three color triplets
is
\begin{equation}
  \theta_{\text{color}} = \tfrac{1}{\sqrt6}\,\epsilon_{abc}\,
     \ket{a}\ket{b}\ket{c},
  \label{eq:colorsinglet}
\end{equation}
with $\epsilon_{abc}$ the $SU(3)$ Levi--Civita symbol. This factor is
\emph{totally antisymmetric}. (In the Tessera construction,
Section~\ref{sec:emergent}, this $\epsilon_{abc}$ singlet is realized not as an
imposed index but as the \emph{carried} representation $[1,\omega,\omega^2]$ on
three emergent topological holes; $\omega=e^{2\pi i/3}$.)

\paragraph{The forcing argument.}
For the spatial ground state the three quarks sit in the lowest orbital, which
is \emph{symmetric}. So:
\begin{equation}
  \underbrace{\Psi}_{\text{antisym}} =
  \underbrace{\psi_{\text{space}}}_{\text{sym}} \otimes
  \underbrace{(\phi_{\text{flavor}}\otimes\chi_{\text{spin}})}_{?} \otimes
  \underbrace{\theta_{\text{color}}}_{\text{antisym}}
  \;\Rightarrow\;
  \phi_{\text{flavor}}\otimes\chi_{\text{spin}} \text{ must be \emph{symmetric}.}
\end{equation}
(antisym $\times$ ? $\times$ antisym $\times$ sym $=$ antisym forces ? $=$ sym).
There are two ways to build a symmetric flavor--spin factor:
\begin{enumerate}[label=(\alph*),leftmargin=2.2em]
  \item \textbf{symmetric flavor $\otimes$ symmetric spin} --- both fully
        symmetric. Symmetric spin of three is the quartet, $J=\tfrac32$. This is
        the $\Delta$: e.g.\ $\Delta^{++}=\ket{uuu}$ with
        $J^2=\tfrac{15}{4}$.
  \item \textbf{mixed flavor $\otimes$ mixed spin}, combined into a symmetric
        whole via \eqref{eq:msmssym}. Mixed spin of three is the doublet,
        $J=\tfrac12$. This is the nucleon (proton/neutron), $J^2=\tfrac34$.
\end{enumerate}

\paragraph{Why $\tfrac34$ needs the $2{+}1$ (uud) flavor structure.}
The proton has flavor content $uud$ --- \emph{two} $u$ quarks and \emph{one}
$d$. This $2{+}1$ split is what \emph{selects} the mixed--symmetry spin
combination \eqref{eq:protonstate}. Concretely: the two like quarks (the $uu$
pair) sit in the flavor state symmetric under their exchange, and to make the
total flavor--spin factor symmetric \eqref{eq:msmssym} the $uu$ pair's
\emph{spin} must also be symmetric, i.e.\ in the spin--triplet $S_{uu}=1$. The
full $J=\tfrac12$ proton spin--flavor wavefunction is the SU(6) state
\begin{equation}
  \ket{p\!\uparrow} = \tfrac{1}{\sqrt{18}}\big(
      2\ket{u\!\up u\!\up d\!\down}
    - \ket{u\!\up u\!\down d\!\up}
    - \ket{u\!\down u\!\up d\!\up} + \text{(permutations of which quark is }d)\big),
  \label{eq:su6}
\end{equation}
whose spin part, at fixed flavor ordering $uud$, is exactly the
$2\ket{\up\up\down}-\ket{\up\down\up}-\ket{\down\up\up}$ of
\eqref{eq:protonstate}. \textbf{The upshot for the code:} you cannot identify
the $J=\tfrac12$ combination, and hence the value $\tfrac34$, without keeping
the flavor labels and the spin labels \emph{jointly}; if you sum spins blindly
over three indistinguishable holes you land on the symmetric/quartet or the
product mixture, not the mixed--symmetry proton. This is the physics reason the
codebase must read flavor and spin \emph{together}
(Sections~\ref{sec:composite}--\ref{sec:totalspace}).

% =============================================================================
\section{A primer on discrete (lattice) geometry}
\label{sec:lattice}
% =============================================================================

The Tessera engine does not work with smooth fields on a manifold; it works
with a \emph{simplicial complex} --- a triangulation --- and discrete
differential forms living on it. This section fixes that vocabulary.

\paragraph{Simplicial complexes.}
A $k$--\emph{simplex} is the convex hull of $k+1$ affinely independent points:
a vertex ($0$--simplex), edge ($1$--simplex), triangle ($2$--simplex),
tetrahedron ($3$--simplex), pentatope/$4$--simplex, and so on. A
\emph{simplicial complex} $\mathcal K$ is a collection of simplices glued so
that every face of a simplex is in $\mathcal K$ and any two simplices meet in a
common face. The proton host in Section~\ref{sec:emergent} is a triangulated
closed $4$--manifold, a combinatorial $S^4$ built from pentatopes.

\paragraph{Cochains = discrete differential forms.}
A $k$--\emph{cochain} assigns a number (here a complex number) to every oriented
$k$--simplex. The space of $k$--cochains is $C^k$. A $0$--cochain is a function
on vertices; a $1$--cochain is a number per edge (a discrete $1$--form); a
$2$--cochain, a number per triangle; etc. Cochains are the lattice stand--ins
for differential forms $\Omega^k$.

\paragraph{Coboundary $d$ and boundary $\partial$.}
The \emph{coboundary} operator $d:C^k\to C^{k+1}$ is the discrete exterior
derivative: $(d\alpha)(\sigma) = \sum_{\tau\in\partial\sigma} [\sigma:\tau]\,
\alpha(\tau)$, summing the cochain over the oriented faces $\tau$ of $\sigma$
with incidence signs $[\sigma:\tau]=\pm1$. It satisfies $d^2=0$ (the discrete
Stokes/Poincar\'e statement). Its adjoint with respect to the inner product
defined by the metric (below) is the \emph{codifferential}
$\delta:C^k\to C^{k-1}$, $\delta = d^\dagger$. The \emph{Hodge Laplacian} is
\begin{equation}
  \Delta = d\delta + \delta d ,
\end{equation}
and its kernel $\ker\Delta = \mathcal H^k$ is the space of \emph{harmonic}
$k$--cochains, whose dimension is the $k$--th Betti number $b_k$ (a topological
invariant counting $k$--dimensional ``holes''). These harmonic cochains are
exactly the carried register states of Section~\ref{sec:emergent}.

\paragraph{Regge calculus: geometry from edge lengths.}
Regge calculus encodes the geometry not by a smooth metric tensor but by the
set of \emph{squared edge lengths} $\ell_e^2$. Curvature concentrates on the
codimension--$2$ simplices (the ``hinges'': triangles in $4$D). Around a hinge
$h$ the surrounding flat simplices fail to close up by a \emph{deficit angle}
\begin{equation}
  \varepsilon_h = 2\pi - \sum_{\sigma\supset h}\theta_{\sigma,h},
\end{equation}
where $\theta_{\sigma,h}$ is the dihedral angle that simplex $\sigma$ subtends
at the hinge $h$. The deficit $\varepsilon_h$ is the discrete Ricci scalar
curvature smeared on $h$. The (Euclidean) Regge action is
\begin{equation}
  S_{\text{Regge}} = \sum_{\text{hinges } h} A_h\,\varepsilon_h
   \qquad(\text{$A_h$ = the hinge's volume}),
\end{equation}
the lattice version of the Einstein--Hilbert action $\int\!\sqrt{g}\,R$.

\paragraph{Dual (Voronoi/circumcentric) volumes.}
The inner product on cochains needs a volume per cell. Each $k$--simplex
$\sigma$ is paired with its \emph{dual} $(d-k)$--cell, built from circumcenters
(the circumcentric/Voronoi dual). The product of primal and dual volumes,
$V_\sigma V^\star_\sigma$, sets the Hodge star and hence $\delta$, $\Delta$, and
the harmonic inner products. (Concretely the codebase computes these as the
``option A'' circumcentric dual volumes on each \code{Simplex}.)

\paragraph{The Lorentzian (complex) action.}
Spacetime is Lorentzian, not Euclidean: one direction is timelike. On the
lattice this means some squared edge lengths are negative ($\ell_e^2<0$ for
timelike edges), and the dihedral/deficit angles become \emph{boosts}
(hyperbolic angles, $\operatorname{arccosh}$) rather than ordinary rotations.
The clean bookkeeping (Sorkin; Asante--Dittrich) makes the Regge action
\emph{complex}:
\begin{equation}
  S \in \C, \qquad
  \text{Re}\,S \text{ from spacelike hinges}, \quad
  \text{Im}\,S \text{ from timelike (boost) hinges}.
\end{equation}
\textbf{The imaginary part is real physics}, not a numerical artifact: it
constrains degrees of freedom (e.g.\ the overall scale) that the real part
leaves loose. Throughout the engine one therefore works with the full complex
$S$ and objectives like $\sum_e|\partial_e S|^2$, never $\text{Re}\,S$ alone.
This is background for why the host geometry can carry a definite spin structure
at all.

% =============================================================================
\section{The Dirac--K\"ahler framework: forms as spinors}
\label{sec:dk}
% =============================================================================

We now need spin on a lattice. The standard continuum Dirac field needs a frame
(vierbein) at each point and is awkward to discretize; the \emph{Dirac--K\"ahler}
(DK) formulation, by contrast, packages spinors as \emph{inhomogeneous
differential forms} --- exactly the cochains of Section~\ref{sec:lattice} ---
and so discretizes naturally.

\paragraph{The Dirac--K\"ahler equation.}
Collect a form of each degree into one inhomogeneous form
$\Phi = \sum_{k=0}^{d}\phi^{(k)}$, $\phi^{(k)}\in\Omega^k$. Define the operator
\begin{equation}
  D \eqdef d + \delta .
\end{equation}
Because $d^2=\delta^2=0$, it squares to the Laplacian,
\begin{equation}
  D^2 = (d+\delta)^2 = d\delta + \delta d = \Delta ,
\end{equation}
so $D$ is a ``square root of the Laplacian'' --- a Dirac--type operator. The
Dirac--K\"ahler equation $(D + m)\Phi = 0$ is, after a change of variables,
equivalent to several copies of the Dirac equation.

\paragraph{The K\"ahler--Atiyah isomorphism.}
The exterior algebra $\Lambda(\C^d)$ of forms, equipped with the Clifford
(``geometric'') product, is isomorphic as an algebra to a matrix algebra. For
$d=4$:
\begin{equation}
  \boxed{\,\Lambda(\C^4) \;\cong\; M_4(\C), \qquad
          \dim = 2^4 = 16 = \underbrace{4}_{\text{Dirac spinor}}
                              \times \underbrace{4}_{\text{taste}}.\,}
  \label{eq:KA}
\end{equation}
The $16$--dimensional space of inhomogeneous forms factorizes as a $4$--component
Dirac spinor times a $4$--fold ``taste'' (a discretization multiplicity that, in
this codebase, plays the role of the per--hole \emph{flavor} index --- the
``taste'' factor of the DK fiber). The spinor factor carries the spin; the taste
factor carries flavor. This single isomorphism is why one mesh can carry
\emph{both} the spin and the flavor of a quark.

\paragraph{The Clifford generators $\gamma^a$.}
The image of the basis $1$--forms under \eqref{eq:KA} are the gamma matrices
$\gamma^a$ ($a=0,\dots,3$), the Clifford action of a $1$--cochain. They satisfy
\begin{equation}
  \{\gamma^a,\gamma^b\} = 2\,\eta^{ab}\,\id ,
  \label{eq:clifford}
\end{equation}
with $\eta^{ab}=\operatorname{diag}(-,+,+,+)$ the (Lorentzian) signature; in the
code these are \code{DiracKahler.gammas}, of dimension
\code{gammaDimension()}$=4$, and the function \code{clifford\_residual} measures
the deviation of \eqref{eq:clifford} from $0$ --- a certificate that the
generators close.

\paragraph{Spin generators and the structural spin--\textonehalf.}
The Lorentz/rotation generators in the Clifford algebra are
\begin{equation}
  \boxed{\,\Sigma_{ij} = \tfrac14\,[\gamma_i,\gamma_j].\,}
  \label{eq:Sigma}
\end{equation}
For two distinct \emph{spatial} indices $i\neq j$ (with
$\gamma_i^2=\gamma_j^2=+\id$ and $\gamma_i\gamma_j=-\gamma_j\gamma_i$),
$\Sigma_{ij}=\halff\gamma_i\gamma_j$ and
$\Sigma_{ij}^2 = \tfrac14\gamma_i\gamma_j\gamma_i\gamma_j
= -\tfrac14\gamma_i^2\gamma_j^2 = -\tfrac14\id$, so $\Sigma_{ij}$ has
eigenvalues $\pm\tfrac{i}{2}$, i.e.\ the \emph{spin component}
$S_{ij}=-i\Sigma_{ij}$ has eigenvalues $\pm\tfrac12$. Forming the spin Casimir
over the three spatial planes,
\begin{equation}
  S^2_{\text{DK}} = \sum_{i<j\,\in\{1,2,3\}} S_{ij}^2
                  = 3\cdot\tfrac14\,\id = \tfrac34\,\id .
  \label{eq:dkcasimir}
\end{equation}
\emph{Every Dirac--K\"ahler field is structurally spin--\textonehalf}: the
Casimir is $\tfrac34$, the same fingerprint as a single $\C^2$
(Section~\ref{sec:onespin}). A scalar would give $0$ and a vector $0,\pm1$; only
a spinor gives the uniform $\pm\tfrac12$ eigenvalues. This is what
\code{dk\_spin\_readout.py} certifies on the bare mesh.

\paragraph{Structural vs.\ composite spin --- do not confuse them.}
There are now \emph{two} different ``$\tfrac34$''s, and keeping them apart is
essential:
\begin{itemize}[leftmargin=1.6em]
  \item the \textbf{structural} spin--\textonehalf{} of \eqref{eq:dkcasimir} is a
        property of \emph{one} DK fiber --- one quark --- guaranteed by the
        algebra. It is the lattice analog of \eqref{eq:onecasimir}.
  \item the \textbf{composite} total spin of Section~\ref{sec:threespin} is a
        property of \emph{three} quarks bound together, and equals $\tfrac34$
        only for the entangled mixed--symmetry combination
        \eqref{eq:protonstate}; for a product it is $\tfrac74$, for the symmetric
        state $\tfrac{15}{4}$.
\end{itemize}
The code reaches the first $\tfrac34$ easily (it is structural); the hard,
partly--open problem is the \emph{second} $\tfrac34$.

% =============================================================================
\section{The emergent proton in the Tessera codebase}
\label{sec:emergent}
% =============================================================================

We can now describe how a proton appears in the engine, and where spin enters.

\paragraph{The host and the optimizer.}
The computation lives on a closed $4$--manifold host (a combinatorial $S^4$ of
pentatopes). A \code{MultiCobordism} optimizer evolves the geometry by gated
local moves (Pachner moves and surgical conings), keeping only moves that
improve the stationary--action objective
\begin{equation}
  F = \norm{\nabla S_{\text{Regge}}}^2 + \Gamma\cdot r_U ,
\end{equation}
which \emph{extremizes} the (complex) action ($\delta S=0$, never minimizes its
magnitude) while $r_U$ pins the carried register. The topology is
\emph{emergent}: holes are not hand--placed but grown under relaxation and kept
only if they lower $F$.

\paragraph{The color register as emergent holes.}
The proton's three colors appear as three topological holes of the host. On a
$4$--manifold the relevant homology is degree $3$: the color register is the
$b_3$ harmonic sector, read at cochain degree $k=3$. Removing $n$ disjoint top
cells produces $b_3 = n-1$ independent holes; three holes is the proton's color
register. The carried register state is the harmonic $3$--cochain
\begin{equation}
  \psi = \texttt{EigenstateSynthesis(st, 3).carriedRepresentative},
\end{equation}
a least--action representative over the holes that carries the color
representation $[1,\omega,\omega^2]$ with $\omega=e^{2\pi i/3}$ --- the discrete
realization of the antisymmetric $\epsilon_{abc}$ singlet
\eqref{eq:colorsinglet}. The carried--singlet condition is certified by the
register residual $r_{\text{state}}\to0$; the singlet $[1,1,1]$ is \emph{not}
carried (residual $\sim 100$), which is exactly confinement.

\paragraph{Per--hole spinor extraction.}
A simplicial complex has no global frame, so spins at different holes are not
\emph{a priori} comparable. The extraction
(\code{dk\_composite\_spin.py}, \code{dk\_spin\_wilson.py}) proceeds:
\begin{enumerate}[leftmargin=1.8em]
  \item \textbf{Cell frame.} For each hole's dual cell $c$, build a frame $F(c)$
        --- the \code{cell\_frame}: centered principal axes (inertia--tensor
        eigenvectors) of the cell's \emph{dual--edge} vectors
        (circumcenter--to--neighboring--circumcenter). Because the dual--edge
        \emph{lengths} carry the neighbors' sizes, the frame is non--degenerate
        even when the cell's own vertices are near--symmetric, and it depends
        only on the vertex \emph{set} (relabel--invariant, gauge--covariant
        $F\to RF$).
  \item \textbf{Spin transport.} The \code{facet\_transport} across a shared
        facet is the $\mathrm{Spin}(4)$ rotation aligning the two frames, lifted
        from $SO(4)$ by \code{rotation\_to\_spin}; \code{wilson\_line}$(i,j)$
        composes these along a BFS dual path --- the discrete
        \emph{spin--connection Wilson line} (open--path holonomy) that makes the
        comparison frame--independent. A rotation by $\epsilon$ produces spinor
        eigenphases $\pm\epsilon/2$, the $\mathrm{Spin}$ double--cover signature
        checked by \code{is\_double\_cover}.
  \item \textbf{Spinor.} \code{emergent\_spinor} reads the per--hole spinor from
        $\psi$ by least--squaring the cell's tetrahedral trivector minors against
        $\psi$ in $F(c)$ and mapping the resulting $3$--form into the Clifford
        algebra, $\Phi=\sum_t\omega_t\,\gamma_i\gamma_j\gamma_k$, then
        $s=\Phi\cdot[1,0,0,0]$.
\end{enumerate}
The validation gates (GAUGE invariance, RELABEL invariance) are run
\emph{post--hoc} in the test module and \code{\_\_main\_\_}, never as a loop
condition. Reference modules:
\code{examples/cobordism/emergent\_optimizer.py},
\code{examples/cobordism/dk\_composite\_spin.py},
\code{examples/cobordism/dk\_joint\_spin.py},
\code{examples/cobordism/dk\_spin\_readout.py}.

% =============================================================================
\section{The composite--spin readout problem (the heart)}
\label{sec:composite}
% =============================================================================

With per--hole spinors in hand, the obvious move is to combine them into the
total $J^2$ and check whether the emergent bound state is a proton ($\tfrac34$)
or a $\Delta$ ($\tfrac{15}{4}$). This is where the physics of
Section~\ref{sec:threespin} returns with teeth.

\paragraph{The product read and its provable floor.}
\code{composite\_j2} transports the three holes' spinors to a common frame via
the Wilson lines, forms the \emph{product} state
$\ket{s_1}\otimes\ket{s_2}\otimes\ket{s_3}$, and reads $J^2$. Equivalently, each
hole is reduced to a single--qubit Bloch vector $\expval{\bm S}_h$
(Section~\ref{sec:onespin}). But a product of three spin--\textonehalf{} Bloch
vectors is exactly the object of \eqref{eq:product74}: its $J^2$ is the
mixture value, floored \emph{above} the proton's $\tfrac34$ and unable to reach
it. The code observes precisely this floor: the readout robustly certifies a
\emph{three--spin--\textonehalf{} (baryon)} state --- $|\expval{\bm S}|=\tfrac12$
per hole --- but $J^2$ sits at the $n=3$ mixture ($\sim\tfrac94$ in the joint
estimators), never $\tfrac34$. \textbf{This is not a bug, a convergence
failure, or a pairing problem}: it is the theorem of
Section~\ref{sec:threespin} --- a product cannot represent an entangled
mixed--symmetry state.

\paragraph{The pair--creation build (correcting the two--quark artifact).}
An earlier attempt read the joint state off a \emph{static}, hand--made color
singlet fed only \emph{two} quarks, and found the per--hole flavor index
\emph{not} independent --- an artifact of underfeeding a state that
intrinsically needs three quarks. The correct construction \emph{builds} the
proton rather than reading off a hand--made singlet. One co--optimized
\code{MultiCobordism} runs a \emph{simultaneous} pair--creation event,
\begin{equation}
  3 \text{ neutral } q\bar q \text{ pairs } (\Sigma=0)
  \;\longrightarrow\; [\,\text{proton},\ \text{antiproton}\,],
\end{equation}
with the proton and antiproton legs co--optimized in \emph{one} system (charge
conserved end to end). The proton's three quarks are the three emergent holes of
the proton output block, accessed read--only via
\code{MultiCobordism.outputs[0].verts} with the relaxed metric copied onto the
sub--complex. This build delivers, and validates, two things:
\begin{itemize}[leftmargin=1.6em]
  \item it \textbf{carries the color singlet} ($r_{\text{state}}\to0$ on
        converged seeds), and
  \item it gives \textbf{independent per--hole flavor} (the DK taste/charge is
        distinguishable across the three quark holes, relative spread
        $\sim0.2$--$0.5$) --- exactly the three--quark structure the two--quark
        read lacked.
\end{itemize}

\paragraph{And yet $J^2$ still sits at the mixture.}
Even on this correct build, with carried singlet and independent flavor, the
composite $J^2$ remains at the indefinite mixture ($\sim\tfrac94$), not
$\tfrac34$:
\begin{center}
\begin{tabular}{lcc}
\toprule
flavor handling & joint $J^2$ & per--hole $J^2$ \\
\midrule
product           & $\approx1.5$ & $\approx1.8$ \\
$2{+}1$ clustering & $\approx2.2$ & $\approx2.8$ \\
\bottomrule
\end{tabular}
\end{center}
The reason is now familiar: \emph{every available readout reduces each hole to a
single--qubit Bloch vector}, so the three--quark spin state it reconstructs is a
\emph{product}, and a product provably cannot be the entangled $\tfrac34$ (recall
\eqref{eq:product74}: even the product $\ket{uud}$ is $\tfrac74$). The build
\emph{creates} the entanglement; the per--hole readout \emph{discards} it. Even
charge--clustering the flavors into a $2{+}1$ ($u/d$) pairing and symmetrizing
the $uu$ spins does not help --- $\ket{1,1}\ket{\tfrac12,-\tfrac12}$ is still
$\tfrac74$, not the pure $J=\tfrac12$ eigenstate. The validated regression guard
(in \code{test\_dk\_joint\_spin.py}) confirms the \emph{measuring stick} is
correct: fed the clean states by hand, the $J^2$ operator returns
$\tfrac34$ (proton eigenstate), $\tfrac{15}{4}$ ($\Delta$, $\ket{uuu}$), and
$\tfrac74$ (product $\ket{uud}$). The instrument is right; the per--hole
\emph{sampling} of the emergent state is what loses the entanglement.

% =============================================================================
\section{The total--space spin operator (the resolution being built)}
\label{sec:totalspace}
% =============================================================================

The diagnosis dictates the cure. To recover $\tfrac34$ the spin operator must
act on the \emph{whole} carried representative at once --- the entangled
three--quark state as one object --- rather than on a product of per--hole
spinors. This is the \emph{total--space spin operator}
(scoped in the codebase as the open follow--up; companion
\code{docs/theory/cobordism/proton-spin/cartan\_weyl\_gluon.tex}).

\paragraph{The operator.}
Let $P_h$ be the projector onto the $h$--th hole's degrees of freedom inside the
single joint carried field $\psi$, and let $\Sigma_a$ ($a$ a spatial rotation
plane) be the structural spin generator \eqref{eq:Sigma} acting on the DK fiber.
The total spin component is $J_a = \sum_h P_h\,\Sigma_a\,P_h$, and the Casimir is
\begin{equation}
  \boxed{\;J^2 = \sum_a \Big(\sum_h P_h\,\Sigma_a\,P_h\Big)^2
              = \sum_a \sum_{h,h'} P_h\Sigma_a P_h\,P_{h'}\Sigma_a P_{h'}\;}
  \label{eq:totalj2}
\end{equation}
acting on the \emph{joint} field. The \emph{diagonal} terms $h=h'$ reproduce the
per--hole pieces (the product floor); the crucial \emph{cross--hole} terms
$h\neq h'$ are precisely the entanglement carriers that the product read threw
away --- they are the lattice analog of the swap operators
$P_{12},P_{13},P_{23}$ in \eqref{eq:j2swap} that produced the cancellation
$2-1-1=0$ giving $\tfrac34$. Keeping these cross terms is the whole point.

\paragraph{The core difficulty: fibers vs.\ cells.}
The $16$--dimensional form fiber on which $\Sigma_a$ acts is defined
\emph{point--wise}, but the cochain components of $\psi$ live on simplices of
\emph{different degrees} (vertices, edges, triangles, tetrahedra) that do not
sit at the same point and do not align point--by--point on a generic simplicial
complex. So one cannot simply ``apply $\Sigma_a$'' to $\psi$. The assembly
needs a map from the cells to the full DK total space --- a
\emph{Whitney}/K\"ahler--Atiyah interpolation that lifts the degree--stratified
cochain $\psi$ into the $16$--component DK field where $\Sigma_a$ is defined,
applies \eqref{eq:totalj2}, and projects back. Building this fiber$\leftrightarrow$cells
map is the open engineering task (issue \#495, the readout follow--up to the
build of \#489), and it is the last rung between the validated $\tfrac74$/$\tfrac{15}{4}$/$\tfrac34$
\emph{instrument} and a $\tfrac34$ read \emph{off the emergent proton}.

% =============================================================================
\section{Summary}
\label{sec:summary}
% =============================================================================

\paragraph{The numbers, all in one place.}
\begin{center}
\renewcommand{\arraystretch}{1.3}
\begin{tabular}{l c l}
\toprule
state & $J^2$ & meaning \\
\midrule
single spin--\textonehalf{} (one $\C^2$) & $\tfrac34$ & structural spin, Eq.\eqref{eq:onecasimir} \\
DK fiber (one quark) & $\tfrac34$ & structural, Clifford $\Sigma_{ij}$, Eq.\eqref{eq:dkcasimir} \\
proton, entangled mixed--symmetry $2\ket{uud}{-}\ket{udu}{-}\ket{duu}$ & $\tfrac34$ & composite $J=\tfrac12$, Eq.\eqref{eq:protonstate} \\
bare product $\ket{uud}$ (per--hole Bloch read) & $\tfrac74$ & the mixture floor, Eq.\eqref{eq:product74} \\
three--quark mixture floor (current readout) & $\sim\tfrac94$ & product of per--hole qubits \\
$\Delta$, symmetric $\ket{uuu}$ & $\tfrac{15}{4}$ & composite $J=\tfrac32$ quartet \\
\bottomrule
\end{tabular}
\end{center}

\paragraph{Validated vs.\ open.}
\begin{itemize}[leftmargin=1.6em]
  \item \textbf{Validated.} Structural spin--\textonehalf{} of the DK fiber
        (\code{dk\_spin\_readout.py}); per--hole spinor extraction with
        gauge/relabel invariance via the dual--edge frame and spin--connection
        Wilson line (\code{dk\_composite\_spin.py}); the emergent color singlet
        $[1,\omega,\omega^2]$ carried with $r_{\text{state}}\to0$; the
        pair--creation build delivering a carried singlet \emph{and} independent
        per--hole flavor (\code{dk\_joint\_spin.py}); and the $J^2$ \emph{instrument}
        itself, which returns $\tfrac34,\ \tfrac74,\ \tfrac{15}{4}$ on
        hand--fed clean states (regression guard \code{test\_dk\_joint\_spin.py}).
  \item \textbf{Open.} Reading $\tfrac34$ off the emergent proton, which requires
        the total--space spin operator \eqref{eq:totalj2} acting on the joint
        carried field --- blocked on the fiber$\leftrightarrow$cells
        (Whitney/K\"ahler--Atiyah) assembly. Tracked as issue \#495 / \#477.
\end{itemize}

\paragraph{The one sentence to remember.}
The proton's spin is an \emph{entangled} three--body quantity; a readout that
reduces each quark to its own Bloch vector lands on the product mixture
($\tfrac74$ per ket, $\sim\tfrac94$ for three), and only an operator acting on
the whole carried representative at once --- keeping the cross--hole terms ---
can recover $\tfrac34$.

\paragraph{Design notes.}
\code{docs/design/joint\_proton\_spin\_findings.md} (the build and the validated
measuring stick, \#489) and \code{docs/theory/cobordism/proton-spin/cartan\_weyl\_gluon.tex}
(the Cartan/Weyl research direction toward it, \#495/\#477).

% =============================================================================
\begin{thebibliography}{9}
\addcontentsline{toc}{section}{References}

\bibitem{becherjoos} E.~Becher and H.~Joos,
  \emph{The Dirac--K\"ahler Equation and Fermions on the Lattice},
  Z.\ Phys.\ C \textbf{15}, 343 (1982).

\bibitem{graf} W.~Graf,
  \emph{Differential forms as spinors},
  Ann.\ Inst.\ Henri Poincar\'e A \textbf{29}, 85 (1978).

\bibitem{sakurai} J.~J.~Sakurai and J.~Napolitano,
  \emph{Modern Quantum Mechanics}, 2nd ed., Addison--Wesley (2011)
  --- addition of angular momenta, Clebsch--Gordan coefficients (Ch.~3).

\bibitem{regge} T.~Regge,
  \emph{General relativity without coordinates},
  Nuovo Cimento \textbf{19}, 558 (1961).

\bibitem{hamber} H.~W.~Hamber,
  \emph{Quantum Gravitation: The Feynman Path Integral Approach},
  Springer (2009) --- Regge calculus and lattice quantum gravity.

\bibitem{sorkin} R.~D.~Sorkin,
  \emph{Lorentzian angles and trigonometry including lightlike vectors},
  arXiv:1908.10022 (2019); and S.~K.~Asante, B.~Dittrich, H.~M.~Haggard,
  \emph{The degrees of freedom of area Regge calculus}, for the complex
  Lorentzian action conventions.

\bibitem{halzenmartin} F.~Halzen and A.~D.~Martin,
  \emph{Quarks and Leptons: An Introductory Course in Modern Particle Physics},
  Wiley (1984) --- $SU(6)$ baryon wavefunctions, the proton/$\Delta$
  spin--flavor structure (Ch.~2--5).

\bibitem{griffiths} D.~Griffiths,
  \emph{Introduction to Elementary Particles}, 2nd ed., Wiley--VCH (2008)
  --- a gentler treatment of the same baryon wavefunctions and the color singlet.

\end{thebibliography}

\end{document}
